\chapter{Špecifikácia požiadaviek}
\label{ch:problem_statement}

Táto kapitola sa venuje analýze súčasného stavu procesov obstarávania
v~študentských radách, identifikácii ich hlavných nedostatkov
a~formulácii požiadaviek na navrhované riešenie. Cieľom je vytvoriť
presné východisko pre návrh systému, ktorý bude tieto nedostatky
adresovať prostredníctvom decentralizovaných technológií. Požiadavky
vychádzajú z~poznatkov získaných v~analytickej časti práce
(kapitola~\ref{ch:newchapter}) a~z~praktických skúseností s~fungovaním
študentských organizácií.

%%%%%%%%%%%%%%%%%%%%%%%%%%%%%%%%%%%%%%%%%%%%%%%%%%%%%%
\section{Analýza súčasného stavu}
\label{sec:current_state}

Študentské rady vo vzdelávacích inštitúciách pravidelne uskutočňujú
menšie nákupy pre potreby svojej činnosti -- materiály na podujatia,
kancelárske potreby, občerstvenie, propagačné predmety a~podobne.
Hoci jednotlivé sumy sú spravidla malé, ich kumulatívny objem počas
akademického roka je značný a~vyžaduje systematický prístup
k~evidencii a~rozhodovaniu.

V~súčasnosti tieto procesy typicky prebiehajú nasledovne:
\begin{itemize}
    \item \textbf{Evidencia financií} je vedená v~tabuľkových
    procesoroch (Microsoft Excel, Google Sheets). Záznamy o~nákupoch,
    rozpočtoch a~výdavkoch sú uložené lokálne alebo v~cloudovom
    úložisku, pričom prístup k~nim má obmedzený okruh osôb.

    \item \textbf{Komunikácia s~dodávateľmi} prebieha neštruktúrovane
    prostredníctvom e-mailov, telefonátov alebo osobných kontaktov.
    Neexistuje jednotný systém na zber a~porovnávanie ponúk.

    \item \textbf{Rozhodovanie o~nákupoch} sa uskutočňuje na stretnutiach
    rady, kde sa diskutuje o~prijatých ponukách a~hlasuje sa o~výbere
    dodávateľa. Proces je manuálny a~závisí od fyzickej prítomnosti
    členov.

    \item \textbf{Dokumentácia} je fragmentovaná -- ponuky, faktúry
    a~zápisnice sa uchovávajú v~rôznych formátoch a~na rôznych
    miestach, bez centrálneho systému na ich správu a~overenie.
\end{itemize}

Takýto spôsob fungovania je pre menšie organizácie bežný a~pri
nízkych objemoch transakcií môže byť postačujúci. S~rastúcim počtom
nákupov a~zvyšujúcimi sa požiadavkami na transparentnosť sa však
stáva zdrojom viacerých problémov, ktoré sú identifikované
v~nasledujúcej sekcii.

%%%%%%%%%%%%%%%%%%%%%%%%%%%%%%%%%%%%%%%%%%%%%%%%%%%%%%
\section{Identifikácia problémov}
\label{sec:problem_identification}

Na základe analýzy súčasného stavu bolo identifikovaných päť hlavných
problémov, ktoré navrhované riešenie musí adresovať.

\textbf{P1: Nedostatok transparentnosti.}
Študenti, v~mene ktorých rada nakladá s~finančnými prostriedkami,
spravidla nemajú prístup k~informáciám o~tom, aké nákupy boli
uskutočnené, aké ponuky boli zvažované a~na základe akých kritérií
bol vybraný konkrétny dodávateľ. Neexistuje verejne dostupný
a~overiteľný záznam o~priebehu obstarávacích procesov. Táto
netransparentnosť môže viesť k~nedôvere študentskej komunity voči
rozhodnutiam rady.

\textbf{P2: Riziko manipulácie s~dátami.}
Záznamy vedené v~tabuľkových procesoroch môžu byť kedykoľvek
upravené bez zanechania stopy o~zmene. Neexistuje mechanizmus, ktorý
by zabraňoval spätnej úprave údajov o~ponukách, cenách alebo
výsledkoch hlasovania. Aj keď úmyselná manipulácia nemusí byť
bežná, samotná možnosť jej uskutočnenia oslabuje dôveryhodnosť
celého procesu.

\textbf{P3: Neefektívne rozhodovanie.}
Rozhodovanie o~výbere dodávateľa vyžaduje zorganizovanie stretnutia
rady, na ktorom sa členovia oboznámia s~ponukami a~hlasujú. Tento
proces je časovo náročný a~závisí od dostupnosti členov. Pri menších
nákupoch je administratívna réžia spojená s~rozhodovaním
neprimeraná hodnote obstarávania.

\textbf{P4: Chýbajúca integrita dokumentov.}
Ponuky dodávateľov, technické špecifikácie a~ďalšie dokumenty
existujú ako bežné súbory, ktoré je možné po ich odoslaní zmeniť
alebo nahradiť. Neexistuje mechanizmus na overenie, že dokument
predložený dodávateľom je totožný s~dokumentom, na základe ktorého
rada rozhodovala. Táto skutočnosť komplikuje prípadný spätný audit.

\textbf{P5: Absencia automatizovanej evidencie a~vyhodnotenia.}
Celý proces obstarávania -- od zverejnenia výzvy, cez zbieranie
ponúk, hlasovanie až po vyhlásenie výsledku -- je spravovaný
manuálne. Sledovanie stavov tendrov, počítanie hlasov a~uchovávanie
záznamov závisí od ľudského faktora, čo zvyšuje riziko
administratívnych chýb a~oneskorení.

Uvedené problémy nie sú unikátne pre jednu konkrétnu organizáciu,
ale sú charakteristické pre širší okruh študentských rád
a~podobných samosprávnych orgánov, kde sa rozhoduje o~nakladaní
s~kolektívnymi prostriedkami.

%%%%%%%%%%%%%%%%%%%%%%%%%%%%%%%%%%%%%%%%%%%%%%%%%%%%%%
\section{Funkčné požiadavky}
\label{sec:functional_requirements}

Funkčné požiadavky definujú špecifické správanie a~funkcie systému,
ktoré musia byť implementované, aby sa dosiahli stanovené ciele.
Požiadavky sú rozdelené podľa troch hlavných aktérov systému:

\begin{itemize}
    \item \textbf{Člen študentskej rady} -- osoba zodpovedná za
    vytváranie a~správu tendrov, posudzovanie žiadostí dodávateľov,
    vyhodnocovanie ponúk a~hlasovanie.

    \item \textbf{Dodávateľ} -- fyzická alebo právnická osoba, ktorá
    sa registruje do systému a~po schválení predkladá ponuky do
    zverejnených tendrov.

    \item \textbf{Verejnosť (pozorovateľ)} -- akýkoľvek používateľ,
    ktorý má možnosť prezerať transparentné údaje o~prebiehajúcich
    a~ukončených tendroch bez potreby registrácie či pripojenia
    peňaženky.
\end{itemize}

%%%%%%%%%%%%%%
\subsection{Požiadavky pre člena študentskej rady}
\label{subsec:fr_council}

\begin{description}
    \item[FR1: Vytvorenie tendra.] Člen rady musí mať možnosť vytvoriť
    nový tender zadaním nasledujúcich parametrov: názov, textový opis
    predmetu obstarávania, maximálny rozpočet, kategória obstarávania
    a~voliteľne sprievodná dokumentácia vo formáte PDF. Systém musí
    umožňovať priloženie dokumentácie nahratím súboru do
    decentralizovaného úložiska IPFS a~uložením jeho obsahového
    identifikátora CID spolu s~metadátami tendra na blockchain.

    \item[FR2: Dvojkrokové publikovanie.] Vytvorenie tendra musí byť
    rozdelené na dva kroky: najprv vytvorenie konceptu (draft), kde je
    možné overiť a~prípadne upraviť údaje vrátane CID dokumentácie,
    a~následne publikovanie tendra so zadaním lehoty na predkladanie
    ponúk. Tento prístup znižuje riziko chybného zverejnenia.

    \item[FR3: Správa životného cyklu tendra.] Člen rady, ktorý
    tender vytvoril, musí mať možnosť spustiť hlasovanie, finalizovať
    tender po ukončení hlasovania, označiť tender ako splnený
    a~zrušiť tender v~počiatočných fázach procesu.

    \item[FR4: Hlasovanie o~ponukách.] Každý člen rady musí mať
    možnosť odovzdať hlas práve jednej ponuke v~rámci konkrétneho
    tendra. Hlasovanie musí prebiehať on-chain, čím sa zaručí
    transparentnosť a~auditovateľnosť výsledkov. Systém musí
    zabrániť duplicitnému hlasovaniu tej istej osoby.

    \item[FR5: Posudzovanie žiadostí dodávateľov.] Člen rady musí
    mať možnosť schváliť alebo zamietnuť žiadosť externého
    subjektu o~registráciu ako dodávateľ. Schválením sa žiadateľovi
    udelí oprávnenie predkladať ponuky.

    \item[FR6: Prezeranie ponúk a~výsledkov.] Člen rady musí mať
    prístup k~zoznamu všetkých predložených ponúk ku konkrétnemu
    tendru vrátane navrhovanej ceny, lehoty dodania a~opisu, ako aj
    k~výsledkom hlasovania po ukončení tendra.
\end{description}

%%%%%%%%%%%%%%
\subsection{Požiadavky pre dodávateľa}
\label{subsec:fr_vendor}

\begin{description}
    \item[FR7: Registrácia do systému.] Dodávateľ sa musí mať možnosť
    zaregistrovať prostredníctvom vyplnenia žiadosti obsahujúcej názov
    spoločnosti, kontaktné údaje a~stručný opis činnosti. Žiadosť
    musí byť zaznamenaná na blockchaine a~podliehať schváleniu
    členom študentskej rady.

    \item[FR8: Prezeranie otvorených tendrov.] Registrovaný dodávateľ
    musí mať možnosť prezerať zoznam všetkých aktívnych tendrov,
    ich detaily vrátane rozpočtu, kategórie, lehoty uzávierky
    a~prípadnej dokumentácie uloženej v~IPFS.

    \item[FR9: Predloženie ponuky.] Registrovaný dodávateľ musí mať
    možnosť predložiť ponuku ku konkrétnemu tendru zadaním cenovej
    ponuky, lehoty dodania a~textového opisu riešenia. Systém musí
    validovať, že cena nepresahuje maximálny rozpočet tendra, že
    tender je v~stave prijímajúcom ponuky a~že lehota uzávierky
    ešte neuplynula.

    \item[FR10: Overenie ponuky na blockchaine.] Po podaní ponuky musí
    systém umožniť overenie jej zaznamenania na blockchaine
    prostredníctvom transakčného hashu a~udalostí emitovaných smart
    kontraktom.
\end{description}

%%%%%%%%%%%%%%
\subsection{Požiadavky pre verejnosť}
\label{subsec:fr_public}

\begin{description}
    \item[FR11: Prezeranie tendrov bez registrácie.] Systém musí
    umožňovať akémukoľvek používateľovi prezerať zoznam všetkých
    tendrov, ich základné metadáta a~aktuálny stav bez potreby
    pripojenia kryptografickej peňaženky alebo registrácie. Tým
    sa zabezpečuje transparentnosť procesu pre celú študentskú
    komunitu.

    \item[FR12: Overenie výsledkov.] Verejnosť musí mať možnosť
    overiť výsledky hlasovania a~víťaznú ponuku prostredníctvom
    verejne dostupných údajov na blockchaine, prípadne
    prostredníctvom blockchain exploreru.

    \item[FR13: Prístup k~dokumentácii.] Systém musí umožňovať
    prístup k~dokumentácii tendra prostredníctvom IPFS, kde si
    používateľ dokument môže zobraziť alebo stiahnuť pomocou CID
    uloženého v~kontrakte.
\end{description}

%%%%%%%%%%%%%%%%%%%%%%%%%%%%%%%%%%%%%%%%%%%%%%%%%%%%%%
\section{Nefunkčné požiadavky}
\label{sec:nonfunctional_requirements}

Nefunkčné požiadavky definujú kvalitatívne atribúty systému, ktoré
sú kritické pre jeho úspešné fungovanie. Sú rozdelené do piatich
kategórií.

\subsection{Bezpečnosť}
\begin{description}
    \item[NFR1:] Všetky transakcie musia byť podpísané privátnym
    kľúčom používateľa prostredníctvom kryptografickej peňaženky,
    čím sa zabezpečuje autenticita každej operácie.

    \item[NFR2:] Smart kontrakt musí byť otestovaný automatizovanými
    testami a~manuálne skontrolovaný pred nasadením na sieť.

    \item[NFR3:] Systém musí implementovať kontrolu prístupu založenú
    na rolách (RBAC), aby boli kritické operácie dostupné len
    oprávneným používateľom.

    \item[NFR4:] Vstupné údaje musia byť validované priamo v~smart
    kontrakte, aby sa predišlo ukladaniu neplatných alebo škodlivých
    dát na blockchain.
\end{description}

\subsection{Výkonnosť}
\begin{description}
    \item[NFR5:] Potvrdenie transakcie by malo byť dostupné do
    niekoľkých sekúnd vďaka použitiu siete Polygon, ktorá poskytuje
    výrazne kratšie časy bloku v~porovnaní s~hlavnou sieťou Ethereum.

    \item[NFR6:] Transakčné poplatky (gas fees) by mali byť
    minimálne a~pre bežné operácie by nemali presiahnuť zlomky centu,
    čo je dosiahnuteľné na testovacej sieti Polygon Amoy.

    \item[NFR7:] Operácie čítania dát z~kontraktu (view funkcie) musia
    byť dostupné bez transakčných poplatkov a~s~okamžitou odozvou.
\end{description}

\subsection{Použiteľnosť}
\begin{description}
    \item[NFR8:] Používateľské rozhranie musí byť intuitívne
    a~prehľadné, aby systém mohli používať aj osoby bez predchádzajúcich
    skúseností s~blockchain technológiou.

    \item[NFR9:] Systém musí podporovať pripojenie prostredníctvom
    peňaženky MetaMask, ktorá je najrozšírenejšou peňaženkou
    pre interakciu s~EVM-kompatibilnými sieťami.

    \item[NFR10:] Rozhranie musí byť responzívne a~použiteľné na
    zariadeniach s~rôznymi veľkosťami obrazovky.
\end{description}

\subsection{Spoľahlivosť a dostupnosť}
\begin{description}
    \item[NFR11:] Kľúčové dáta (tendre, ponuky, hlasy) musia byť
    uložené na blockchaine, čím je garantovaná ich perzistencia
    a~odolnosť voči výpadkom centrálnych serverov.

    \item[NFR12:] Dokumenty uložené v~IPFS musia byť dostupné
    prostredníctvom mechanizmu pinning, ktorý zabezpečuje ich
    dlhodobú perzistenciu v~sieti.

    \item[NFR13:] Údaje zaznamenané na blockchaine musia byť nemenné
    -- po vytvorení záznamu nesmie byť možné ho spätne modifikovať
    ani vymazať.
\end{description}

\subsection{Transparentnosť}
\begin{description}
    \item[NFR14:] Všetky kľúčové operácie musia byť overiteľné
    prostredníctvom blockchain exploreru (Polygonscan), kde si
    ktokoľvek môže overiť obsah transakcií a~emitovaných udalostí.

    \item[NFR15:] Systém musí v~reálnom čase zobrazovať aktuálny stav
    každého tendra, aby bolo zrejmé, v~akej fáze sa obstarávací
    proces nachádza.
\end{description}

%%%%%%%%%%%%%%%%%%%%%%%%%%%%%%%%%%%%%%%%%%%%%%%%%%%%%%
\section{Prípady použitia}
\label{sec:use_cases}

Diagram prípadov použitia (Use Case Diagram) znázorňuje interakcie
medzi aktérmi a~systémom. Hlavnými aktérmi sú člen študentskej rady,
dodávateľ a~verejnosť (pozorovateľ).

\begin{figure}[H]
    \centering
    \includegraphics[width=0.8\textwidth]{figures/diagram-use-case.png}
    \caption{Diagram prípadov použitia pre systém mikro-tendrov}
    \label{fig:use_case_diagram}
\end{figure}

Z~diagramu vyplývajú nasledujúce hlavné prípady použitia:
\begin{itemize}
    \item \textbf{UC1: Vytvorenie a~publikovanie tendra} -- člen rady
    vytvorí koncept tendra, prípadne nahrá dokumentáciu do IPFS,
    a~následne tender publikuje s~definovanou lehotou.

    \item \textbf{UC2: Registrácia dodávateľa} -- externý subjekt
    podá žiadosť o~registráciu, ktorú člen rady schváli alebo
    zamietne.

    \item \textbf{UC3: Predloženie ponuky} -- registrovaný dodávateľ
    podá ponuku ku konkrétnemu aktívnemu tendru.

    \item \textbf{UC4: Hlasovanie o~ponukách} -- členovia rady hlasujú
    za preferovanú ponuku v~rámci hlasovacej fázy tendra.

    \item \textbf{UC5: Finalizácia a~uzavretie tendra} -- tvorca
    tendra ukončí hlasovanie, systém určí víťaznú ponuku a~tender
    je administratívne uzavretý.

    \item \textbf{UC6: Prezeranie tendrov} -- akýkoľvek používateľ
    vrátane verejnosti si môže prezerať zoznam tendrov, ich detaily,
    ponuky a~výsledky hlasovania.
\end{itemize}

Jednotlivé prípady použitia pokrývajú celý životný cyklus tendra od
jeho vytvorenia, cez registráciu dodávateľov a~predkladanie ponúk, až
po hlasovanie a~finálne vyhodnotenie. Podrobný návrh procesov
a~interakcií je rozpracovaný v~kapitole~\ref{ch:solution_approach}.
